% =====================================================================
%  冯宇杰 / Yujie Feng — Bilingual Academic Résumé
%  Compile with: xelatex resume.tex   (run twice for refs)
%  Theme: NJU Purple (#7B2D8E) + Blue accent (#2563EB)
% =====================================================================

\documentclass[11pt,a4paper]{article}

% --- Languages & fonts ---
\usepackage{ctex}                 % Chinese support
\usepackage[T1]{fontenc}
\usepackage{lmodern}              % Latin Modern
\usepackage{fontawesome5}         % Icons (envelope, map-marker, etc.)

% --- Page layout ---
\usepackage[a4paper, margin=1.6cm, top=1.4cm, bottom=1.6cm]{geometry}

% --- Colors ---
\usepackage[table]{xcolor}
\definecolor{njupurple}{HTML}{7B2D8E}
\definecolor{njupurpledark}{HTML}{5B1F7A}
\definecolor{njupurplelight}{HTML}{F3EBFA}
\definecolor{njupurpleborder}{HTML}{E5D6F0}
\definecolor{njublue}{HTML}{2563EB}
\definecolor{njubluelight}{HTML}{DBEAFE}
\definecolor{bodygray}{HTML}{374151}
\definecolor{mutedgray}{HTML}{6B7280}
\definecolor{bggray}{HTML}{FAFAFC}

% --- Boxes & tables ---
\usepackage[most]{tcolorbox}
\usepackage{tabularx}
\usepackage{array}

% --- Lists ---
\usepackage{enumitem}
\setlist[itemize]{leftmargin=14pt,nosep,topsep=2pt,parsep=1pt,itemsep=1pt,label={\textcolor{njupurple}{\textbullet}}}

% --- Section formatting ---
\usepackage{titlesec}
\titleformat{\section}
  {\normalsize\bfseries\color{njupurple}}
  {}{0pt}{\MakeUppercase}
  [\vspace{-4pt}\color{njupurple}\hrule height 1.2pt]
\titlespacing*{\section}{0pt}{12pt}{6pt}

\titleformat{\subsection}
  {\small\bfseries\color{njublue}}
  {}{0pt}{}
\titlespacing*{\subsection}{0pt}{8pt}{4pt}

% --- Links ---
\usepackage[hidelinks,colorlinks=true,urlcolor=njublue]{hyperref}
\hypersetup{
  pdftitle={冯宇杰 / Yujie Feng — Resume},
  pdfauthor={冯宇杰 / Yujie Feng},
  pdfsubject={NJU Undergraduate Resume · Chemistry-Biology Interdisciplinary}
}

% --- Page style ---
\pagestyle{empty}
\setlength{\parindent}{0pt}
\setlength{\parskip}{3pt}

% --- Helpers ---
\newcommand{\purplechip}[1]{\colorbox{njupurplelight}{\textcolor{njupurple}{\footnotesize\bfseries\,#1\,}}}
\newcommand{\bluechip}[1]{\colorbox{njubluelight}{\textcolor{njublue}{\footnotesize\bfseries\,#1\,}}}
\newcommand{\cvdate}[1]{\hfill{\color{njublue}\small\bfseries #1}\\[-4pt]}
\newcommand{\cvdateinline}[1]{{\color{njublue}\small\bfseries #1}}
\newcommand{\cvaward}[1]{{\color{njupurple}\bfseries #1}}
\newcommand{\muted}[1]{{\color{mutedgray}\small #1}}

% Compact research card
\newtcolorbox{rscard}{
  colback=white, colframe=njupurpleborder,
  boxrule=0.5pt, leftrule=2.5pt,
  arc=2pt, outer arc=2pt,
  left=8pt, right=8pt, top=6pt, bottom=6pt,
  before skip=4pt, after skip=4pt,
}

\newtcolorbox{compcard}{
  colback=white, colframe=njupurpleborder,
  boxrule=0.5pt, leftrule=2.5pt,
  arc=2pt, outer arc=2pt,
  left=8pt, right=8pt, top=6pt, bottom=6pt,
  before skip=4pt, after skip=4pt,
}

\newtcolorbox{simpleitem}{
  colback=white, colframe=njupurpleborder,
  boxrule=0.5pt, leftrule=2.5pt,
  arc=2pt, outer arc=2pt,
  left=8pt, right=8pt, top=4pt, bottom=4pt,
  before skip=3pt, after skip=3pt,
}

% =====================================================================
\begin{document}

% ------------------------------ HEADER ------------------------------
\begin{tcolorbox}[
  enhanced,
  colback=njupurple, colframe=njupurple, coltext=white,
  arc=0pt, outer arc=0pt, boxrule=0pt,
  left=14pt, right=14pt, top=14pt, bottom=14pt,
  width=\textwidth,
]
\color{white}
{\Huge\bfseries 冯宇杰\;/\;Yujie Feng}\\[6pt]
{\large\bfseries 匡亚明学院\;·\;理科强化班（化生交叉方向）\;·\;2025--2029}\\[-2pt]
{\large Kuang Yaming Honors College\;·\;Chemistry-Biology Interdisciplinary\;·\;Class of 2029}\\[8pt]
{\small\faEnvelope\;\href{mailto:251850043@smail.nju.edu.cn}{251850043@smail.nju.edu.cn}%
\;\quad\faMapMarker*\;南京\,/\,Nanjing, China%
}
\end{tcolorbox}

% ------------------------------ ABOUT ------------------------------
\section*{关于 \quad / \quad About}
南京大学匡亚明学院 2025 级本科生，化生交叉方向。本科阶段从事蛋白质分子动力学、大分子受限输运、E3 泛素连接酶配体开发等研究；熟悉 NAMD / LAMMPS / GROMACS / Schrödinger 等模拟工具与 C / Python 编程。寻找 2029 年推免 / 暑研机会。

\medskip
\itshape\small
First-year undergraduate (Class of 2029) at Kuang Yaming Honors College, Nanjing University, in the Chemistry-Biology Interdisciplinary Track. Research to date spans protein molecular dynamics, confined dynamics of macromolecules, and small-molecule ligand development for E3 ubiquitin ligases. Proficient in NAMD, LAMMPS, GROMACS, Schrödinger, and C / Python. Seeking 推免 (graduate-school recommendation) and summer research opportunities for 2029.
\upshape\normalfont

% ------------------------------ EDUCATION ------------------------------
\section*{教育背景 \quad / \quad Education}

\noindent\textbf{南京大学 · 匡亚明学院\;/\;Nanjing University, Kuang Yaming Honors College} \hfill \cvdateinline{2025.09 -- 2029.06}\\
理科强化班（化生交叉方向）· 理学学士（待授）\;/\;B.S.\ in Science Intensive Program (Chemistry-Biology Interdisciplinary Track)\\
\muted{GPA: --- \quad Rank: ---}\\[2pt]
\textbf{\small 相关课程\,/\,Relevant Coursework:}
\purplechip{细胞生物学\,Cell Biology}\purplechip{生物化学\,Biochemistry}\purplechip{物理化学\,Physical Chemistry}\purplechip{物理生物学\,Physical Biology}\purplechip{机器学习\,ML}

% ------------------------------ RESEARCH ------------------------------
\section*{科研经历 \quad / \quad Research Experience}

% --- Project 1: E3 Ubiquitin Ligase ---
\begin{rscard}
{\bfseries E3 泛素连接酶中新型小分子配体的开发 \;/\; Development of Novel Small-Molecule Ligands for E3 Ubiquitin Ligases} \hfill \cvdateinline{2026.07 -- 至今\,/\,Present}\\[2pt]
{\small\muted{导师：陈志翔研究员 · 中科院生物与化学交叉研究中心 \;/\; Advisor: Researcher Chen Zhixiang, CAS Interdisciplinary Research Center of Biology and Chemistry}}\\[2pt]
\begin{itemize}
  \item 针对 FBXO31 与 FBXO21 这两种蛋白质，结合实验室已设计的小分子。
  \item 利用 Schrödinger、GROMACS、PyMOL 进行分子对接与分子动力学验证。
  \item 参与部分候选小分子的合成过程。
\end{itemize}
{\itshape\small Performed molecular docking and MD validation of lab-designed small molecules targeting FBXO31 and FBXO21 using Schrödinger, GROMACS, and PyMOL; participated in partial synthesis of candidate molecules.}\\[4pt]
\hfill\bluechip{Schrödinger}\bluechip{GROMACS}\bluechip{PyMOL}\purplechip{分子对接\,Docking}\purplechip{MD}
\end{rscard}

% --- Project 2: DPD Macromolecules ---
\begin{rscard}
{\bfseries 大分子受限动力学的理论研究 \;/\; Theoretical Study of Confined Dynamics of Macromolecules} \hfill \cvdateinline{2026.04 -- 2027.03}\\[2pt]
{\small\muted{导师：汪蓉教授 · 化学学院 \;/\; Advisor: Prof.\ Wang Rong, School of Chemistry}}\\[2pt]
\begin{itemize}
  \item 独立主持本科生早期科研课题。
  \item 利用耗散粒子动力学（DPD）等方法，借助 LAMMPS 开源软件及自编 C 语言程序。
  \item 模拟大分子穿孔（translocation）过程。
\end{itemize}
{\itshape\small Principal investigator of an independent undergraduate-initiated research project; employed Dissipative Particle Dynamics (DPD) simulations in LAMMPS and self-written C programs to study macromolecular translocation through pores.}\\[4pt]
\hfill\bluechip{LAMMPS}\bluechip{C}\purplechip{DPD}\purplechip{软物质\,Soft Matter}
\end{rscard}

% --- Project 3: BBL Protein ---
\begin{rscard}
{\bfseries BBL 蛋白质及其突变体的基本结构分析 \;/\; Basic Structural Analysis of BBL Protein and Its Mutants} \hfill \cvdateinline{2025.12 -- 2026.03}\\[2pt]
{\small\muted{导师：董昊教授 · 匡亚明学院 \;/\; Advisor: Prof.\ Dong Hao, Kuang Yaming Honors College}}\\[2pt]
\begin{itemize}
  \item 利用分子动力学（MD）模拟，主要使用 NAMD，并用 VMD 做可视化分析。
  \item 分析特定定点突变引起的氢键、盐桥网络变化。
  \item 阐述其对空间构象与折叠路径的影响。
\end{itemize}
{\itshape\small Performed MD simulations using NAMD; visualized with VMD; analyzed hydrogen-bond and salt-bridge network changes induced by site-directed mutations; elucidated effects on spatial conformation and folding pathways.}\\[4pt]
\hfill\bluechip{NAMD}\bluechip{VMD}\purplechip{MD}\purplechip{蛋白质折叠\,Protein Folding}
\end{rscard}

% ------------------------------ PUBLICATIONS ------------------------------
\section*{论文 \quad / \quad Publications}

\begin{enumerate}[leftmargin=18pt,nosep,label=\arabic*.]
\item 冯宇杰, 董昊. 基于分子动力学模拟的 BBL 蛋白质及其突变体的多角度分析. \textit{南京大学第 29 届基础学科论坛} (2026). \cvaward{三等奖}.\\[2pt]
{\itshape\small Feng Y., Dong H. (2026). \textit{Multi-angle Analysis of BBL Protein and Its Mutants Based on Molecular Dynamics Simulations.} 29th Nanjing University Basic Disciplines Forum. \cvaward{Third Prize}.}
\end{enumerate}

% ------------------------------ SKILLS ------------------------------
\section*{技能 \quad / \quad Skills}

\noindent\begin{tabularx}{\textwidth}{@{}p{3.2cm} X@{}}
\textbf{\small 编程语言 \,/\, Programming} & \purplechip{C}\;\purplechip{Python}\;\purplechip{HTML} \\[3pt]
\textbf{\small 模拟软件 \,/\, Simulation} & \bluechip{GROMACS}\;\bluechip{LAMMPS}\;\bluechip{NAMD}\;\bluechip{Schrödinger}\;\bluechip{Gaussian}\;\bluechip{PyMOL}\;\bluechip{ChemDraw} \\[3pt]
\textbf{\small 研究方法 \,/\, Methods} & \purplechip{分子对接\,Docking}\;\purplechip{MD}\;\purplechip{机器学习\,ML}\;\purplechip{量子化学\,Quantum Chem.}\;\purplechip{物化生实验\,Wet Lab} \\[3pt]
\textbf{\small 数据与系统 \,/\, Data \& Systems} & \bluechip{Origin}\;\bluechip{MATLAB}\;\bluechip{Linux 服务器\,Linux Servers} \\[3pt]
\textbf{\small 语言 \,/\, Languages} & \bluechip{中文（母语）\,Chinese (Native)}\;\bluechip{English（--- 待补\,TBD）} \\
\end{tabularx}

% ------------------------------ AWARDS ------------------------------
\section*{奖项与荣誉 \quad / \quad Awards \& Honors}

\subsection*{奖学金 / 荣誉 \;/\; Scholarships \& Honors}
\muted{[目前无内容 --- TBD: scholarships, honorary titles, etc.]}

\subsection*{竞赛 \;/\; Competitions}
\begin{compcard}
{\bfseries 首届"AI + 合成生物"创新大赛 \;/\; First "AI + Synthetic Biology" Innovation Competition} \hfill \cvdateinline{2026.05 -- 2026.09\,(进行中\,/\,Ongoing)}\\[2pt]
{\small\muted{主办：天津大学合成生物技术全国重点实验室 \;/\; Host: National Key Laboratory of Synthetic Biology Technology, Tianjin University}}\\[2pt]
\begin{itemize}
  \item 利用 mt-grasp 采样方法，结合已有模型进行采样与分子动力学模拟。
  \item 设计出改善溶菌酶（lysozyme）热稳定性的方案。
\end{itemize}
{\itshape\small Employed the mt-grasp sampling method combined with pre-developed models for sampling and MD simulations to design a method improving lysozyme thermal stability.}\\[2pt]
{\small\muted{导师：董昊教授 · 匡亚明学院 \;/\; Advisor: Prof.\ Dong Hao, Kuang Yaming Honors College}}
\end{compcard}

% ------------------------------ INTERESTS ------------------------------
\section*{研究兴趣 \quad / \quad Research Interests}

我特别感兴趣于将计算模拟与实验验证相结合，用于药物与材料的理性设计。两个具体方向：
\itshape\small
\begin{enumerate}[leftmargin=18pt,nosep]
\item 通过干--湿实验闭环加速 hit 发现与 lead 优化；
\item 利用机器学习与分子模拟结合，辅助蛋白质工程与功能材料的改进。
\end{enumerate}
\upshape\normalfont
\medskip
I am particularly interested in integrating computational simulations with experimental validation for rational drug and material design. Two specific directions:
\begin{enumerate}[leftmargin=18pt,nosep]
\item Accelerating hit discovery and lead optimization via closed-loop dry-/wet-lab workflows;
\item Leveraging machine learning combined with molecular simulations to assist protein engineering and the improvement of functional materials.
\end{enumerate}

% ------------------------------ PRACTICE ------------------------------
\section*{社会实践 \quad / \quad Social Practice}

\begin{simpleitem}
\textbf{2026 南星梦想计划 · 江苏省如皋中学团队 队长 · 优秀志愿者 · 服务时长 > 100 小时}\\
{\itshape\small 2026 Nanxing Dream Plan · Rugao High School (Jiangsu) team leader · Outstanding Volunteer · > 100 service hours}
\hfill \cvdateinline{2026}
\end{simpleitem}

\begin{simpleitem}
\textbf{如皋团市委聘"驻宁青春大使"}\\
{\itshape\small Appointed "Youth-in-Nanjing" Ambassador by the Rugao Municipal Committee of the Communist Youth League}
\hfill \cvdateinline{2025.10 -- 2027.11}
\end{simpleitem}

% ------------------------------ STUDENT WORK ------------------------------
\section*{学生工作 \quad / \quad Student Work}

\begin{simpleitem}
南京大学新生学院安邦书院 · 团学联 \textbf{学术科创部 成员}\\
{\itshape\small Member, Academic \& Innovation Dept., Youth League and Student Union Federation, Anbang Academy (School of Freshmen), NJU}
\hfill \cvdateinline{2025.09 -- 2026.06}
\end{simpleitem}

\begin{simpleitem}
南京大学"如皋中学班" \textbf{班长}\\
{\itshape\small \textbf{Class President}, "Rugao High School Class" Cohort, NJU}
\hfill \cvdateinline{2025.09 -- 2026.08}
\end{simpleitem}

\begin{simpleitem}
南京大学招生宣传志愿者协会 · 江苏省如皋中学 \textbf{理事}\\
{\itshape\small \textbf{Director} (Rugao HS constituency), Admissions Promotion Volunteer Association, NJU}
\hfill \cvdateinline{2026.05 -- 2027.05}
\end{simpleitem}

% ------------------------------ REFERENCES ------------------------------
\section*{推荐人 \quad / \quad References}

\muted{[TBD --- 待补充：推荐人姓名 / 职称 / 联系邮箱 \;/\, Recommender name / title / contact email]}

% ------------------------------ FOOTER ------------------------------
\vfill
\begin{center}
\small\color{mutedgray}
\faGithub\;github.com/fragile-0118\quad $\cdot$\quad
\href{https://fragile-0118.github.io/resume/}{fragile-0118.github.io/resume/}\\[2pt]
{\footnotesize © 2026 冯宇杰 / Yujie Feng \quad · \quad 详细版：\href{https://github.com/fragile-0118/resume}{github.com/fragile-0118/resume}}
\end{center}

\end{document}